\begin{frame}{Learning Outcomes}
  \setlength{\parskip}{0pt}
  By the end of this lecture you should be able to:
  \begin{itemize}\footnotesize\itemsep0.1em
    \item \textbf{Formulate} restoration as an inverse problem and \textbf{choose} fidelity or perceptual metrics using the perception-distortion trade-off.
    \item \textbf{Explain} how real-world super-resolution moved to diffusion priors and one-step models, and when generated detail is hallucination.
    \item \textbf{Compare} denoisers and deblurrers for real noise and blur, including training without clean targets and zero-shot posterior sampling.
    \item \textbf{Choose} an inpainting or object-removal model for large masks, shadows and reflections, or video.
    \item \textbf{Design} an industrial anomaly-detection system, from feature memories and student-teacher models to zero-shot inspectors and pixel metrics.
    \item \textbf{Describe} what changes for medical images and how self-configuring, promptable and pathology foundation models respond.
    \item \textbf{Explain} how geospatial foundation models and embedding fields use bands, time and location, and how they compare with a U-Net.
    \item \textbf{Build} a traffic-analytics pipeline and \textbf{explain} bird's-eye-view perception, end-to-end driving, and why open-loop scores mislead.
  \end{itemize}
  \vspace{0.2em}
  {\footnotesize \textbf{Assumed background} (all covered earlier): diffusion, ControlNet, LoRA, denoising autoencoders, LPIPS, CLIP, DINOv2, SAM, vision-language models, Mahalanobis and $k$-NN scores, AUROC, U-Net, YOLO.\par}
\end{frame}
